\needspace{6\baselineskip}
\noindent\textbf{Bex54}\enspace \citet{Bexton1954}.\par
{\small\noindent\textit{Design:} experiment.\quad \textit{Sample:} paid male students.\quad \textit{Measure:} tolerance of reduced sensory variation; cognitive tests.\par
\noindent\textit{Finding.} Prolonged confinement with little sensory variation was found highly aversive; many participants ended their participation early despite the pay, and some cognitive performance declined.\par
\noindent\textit{Codes:} \texttt{under+ seek}.\quad \textit{Bearing:} (a) \textbf{S}, (b) \textbf{S}.\par
\noindent\textit{Rule trace:} (a) S: support {[}under+{]}, challenge {[}--{]}; (b) S: motivation {[}seek{]}, signal {[}--{]}, cost {[}--{]}.\par}
\medskip
\needspace{6\baselineskip}
\noindent\textbf{Ber60}\enspace \citet{Berlyne1960}.\par
{\small\noindent\textit{Design:} theory.\quad \textit{Sample:} none (theory and review).\quad \textit{Measure:} --.\par
\noindent\textit{Finding.} Proposes that organisms seek an intermediate level of arousal and that monotony motivates exploratory and diversive behavior.\par
\noindent\textit{Codes:} \texttt{under+ seek}.\quad \textit{Bearing:} (a) \textbf{S}, (b) \textbf{S}.\par
\noindent\textit{Rule trace:} (a) S: support {[}under+{]}, challenge {[}--{]}; (b) S: motivation {[}seek{]}, signal {[}--{]}, cost {[}--{]}.\par}
\medskip
\needspace{6\baselineskip}
\noindent\textbf{Csi75}\enspace \citet{Csikszentmihalyi1975}.\par
{\small\noindent\textit{Design:} interview-based.\quad \textit{Sample:} adults engaged in demanding activities.\quad \textit{Measure:} interviews and questionnaires.\par
\noindent\textit{Finding.} Introduces the flow model: absorption where challenge and skill are matched, boredom where skill exceeds challenge, and anxiety, rather than boredom, where challenge exceeds skill.\par
\noindent\textit{Codes:} \texttt{under+ fit- overanx}.\quad \textit{Bearing:} (a) \textbf{S/C}, (b) \textbf{N}.\par
\noindent\textit{Rule trace:} (a) S/C: support {[}fit-, under+{]}, challenge {[}overanx{]}; (b) N: motivation {[}--{]}, signal {[}--{]}, cost {[}--{]}.\par}
\medskip
\needspace{6\baselineskip}
\noindent\textbf{Zuc79}\enspace \citet{Zuckerman1979}.\par
{\small\noindent\textit{Design:} psychometric.\quad \textit{Sample:} several adult samples.\quad \textit{Measure:} Sensation Seeking Scale (Boredom Susceptibility subscale).\par
\noindent\textit{Finding.} Presents sensation seeking as a trait with a boredom-susceptibility component: individuals differ in their aversion to repetition and routine.\par
\noindent\textit{Codes:} \texttt{meas person}.\quad \textit{Bearing:} (a) \textbf{N}, (b) \textbf{N}.\par
\noindent\textit{Rule trace:} (a) N: support {[}--{]}, challenge {[}--{]}; (b) N: motivation {[}--{]}, signal {[}--{]}, cost {[}--{]}; accommodated by Cor.~3.4 of EXI-10-a: {[}person{]}.\par}
\medskip
\needspace{6\baselineskip}
\noindent\textbf{OHa81}\enspace \citet{OHanlon1981}.\par
{\small\noindent\textit{Design:} review.\quad \textit{Sample:} none (review of workplace studies).\quad \textit{Measure:} --.\par
\noindent\textit{Finding.} Reviews boredom in repetitive work and proposes that it arises from sustained effort to attend to monotonous tasks.\par
\noindent\textit{Codes:} \texttt{under+ attn}.\quad \textit{Bearing:} (a) \textbf{S}, (b) \textbf{N}.\par
\noindent\textit{Rule trace:} (a) S: support {[}attn, under+{]}, challenge {[}--{]}; (b) N: motivation {[}--{]}, signal {[}--{]}, cost {[}--{]}.\par}
\medskip
\needspace{6\baselineskip}
\noindent\textbf{Far86}\enspace \citet{FarmerSundberg1986}.\par
{\small\noindent\textit{Design:} psychometric.\quad \textit{Sample:} undergraduates.\quad \textit{Measure:} Boredom Proneness Scale (BPS).\par
\noindent\textit{Finding.} Develops a 28-item true--false scale of the disposition to experience boredom, with reported reliability and correlations with measures of depression, hopelessness and loneliness.\par
\noindent\textit{Codes:} \texttt{meas chronic}.\quad \textit{Bearing:} (a) \textbf{N}, (b) \textbf{N}.\par
\noindent\textit{Rule trace:} (a) N: support {[}--{]}, challenge {[}--{]}; (b) N: motivation {[}--{]}, signal {[}--{]}, cost {[}--{]}; accommodated by Cor.~3.4 of EXI-10-a: {[}chronic{]}.\par}
\medskip
\needspace{6\baselineskip}
\noindent\textbf{Csi89}\enspace \citet{CsikszentmihalyiLeFevre1989}.\par
{\small\noindent\textit{Design:} experience sampling.\quad \textit{Sample:} adult workers.\quad \textit{Measure:} experience sampling of challenge, skill and quality of experience.\par
\noindent\textit{Finding.} Moments of high and balanced challenge and skill were reported more often at work than in leisure and were associated with more positive experience.\par
\noindent\textit{Codes:} \texttt{fit-}.\quad \textit{Bearing:} (a) \textbf{S}, (b) \textbf{N}.\par
\noindent\textit{Rule trace:} (a) S: support {[}fit-{]}, challenge {[}--{]}; (b) N: motivation {[}--{]}, signal {[}--{]}, cost {[}--{]}.\par}
\medskip
\needspace{6\baselineskip}
\noindent\textbf{Csi90}\enspace \citet{Csikszentmihalyi1990}.\par
{\small\noindent\textit{Design:} theory.\quad \textit{Sample:} none (synthesis).\quad \textit{Measure:} --.\par
\noindent\textit{Finding.} Synthesizes flow research: optimal experience where challenge matches skill; boredom where skill exceeds challenge; anxiety where challenge exceeds skill.\par
\noindent\textit{Codes:} \texttt{under+ fit- overanx}.\quad \textit{Bearing:} (a) \textbf{S/C}, (b) \textbf{N}.\par
\noindent\textit{Rule trace:} (a) S/C: support {[}fit-, under+{]}, challenge {[}overanx{]}; (b) N: motivation {[}--{]}, signal {[}--{]}, cost {[}--{]}.\par}
\medskip
\needspace{6\baselineskip}
\noindent\textbf{Vod90}\enspace \citet{VodanovichKass1990}.\par
{\small\noindent\textit{Design:} psychometric.\quad \textit{Sample:} undergraduates.\quad \textit{Measure:} BPS (factor analysis).\par
\noindent\textit{Finding.} Finds that the BPS is multidimensional, with factors including the need for external stimulation and difficulty in self-entertainment.\par
\noindent\textit{Codes:} \texttt{meas}.\quad \textit{Bearing:} (a) \textbf{N}, (b) \textbf{N}.\par
\noindent\textit{Rule trace:} (a) N: support {[}--{]}, challenge {[}--{]}; (b) N: motivation {[}--{]}, signal {[}--{]}, cost {[}--{]}.\par}
\medskip
\needspace{6\baselineskip}
\noindent\textbf{Lar91}\enspace \citet{LarsonRichards1991}.\par
{\small\noindent\textit{Design:} experience sampling.\quad \textit{Sample:} middle-school students.\quad \textit{Measure:} experience sampling of boredom.\par
\noindent\textit{Finding.} Boredom was frequent both in and out of school, and students who were often bored in class tended to be often bored elsewhere.\par
\noindent\textit{Codes:} \texttt{person}.\quad \textit{Bearing:} (a) \textbf{N}, (b) \textbf{N}.\par
\noindent\textit{Rule trace:} (a) N: support {[}--{]}, challenge {[}--{]}; (b) N: motivation {[}--{]}, signal {[}--{]}, cost {[}--{]}; accommodated by Cor.~3.4 of EXI-10-a: {[}person{]}.\par}
\medskip
\needspace{6\baselineskip}
\noindent\textbf{San92}\enspace \citet{Sansone1992}.\par
{\small\noindent\textit{Design:} experiment.\quad \textit{Sample:} undergraduates.\quad \textit{Measure:} interest and task-variation strategies on a repetitive task.\par
\noindent\textit{Finding.} When given a reason to persist at a boring task, participants varied the way they performed it, making it more interesting.\par
\noindent\textit{Codes:} \texttt{switch}.\quad \textit{Bearing:} (a) \textbf{N}, (b) \textbf{S}.\par
\noindent\textit{Rule trace:} (a) N: support {[}--{]}, challenge {[}--{]}; (b) S: motivation {[}switch{]}, signal {[}--{]}, cost {[}--{]}.\par}
\medskip
\needspace{6\baselineskip}
\noindent\textbf{Mon96}\enspace \citet{MonetaCsikszentmihalyi1996}.\par
{\small\noindent\textit{Design:} experience sampling.\quad \textit{Sample:} high-school students.\quad \textit{Measure:} experience sampling of challenge, skill and experience.\par
\noindent\textit{Finding.} Perceived challenge, skill and their balance were related to concentration, involvement and other aspects of the quality of experience.\par
\noindent\textit{Codes:} \texttt{fit-}.\quad \textit{Bearing:} (a) \textbf{S}, (b) \textbf{N}.\par
\noindent\textit{Rule trace:} (a) S: support {[}fit-{]}, challenge {[}--{]}; (b) N: motivation {[}--{]}, signal {[}--{]}, cost {[}--{]}.\par}
\medskip
\needspace{6\baselineskip}
\noindent\textbf{Som00}\enspace \citet{SommersVodanovich2000}.\par
{\small\noindent\textit{Design:} correlational.\quad \textit{Sample:} undergraduates.\quad \textit{Measure:} BPS; symptom checklists.\par
\noindent\textit{Finding.} Boredom proneness was associated with reports of psychological and physical-health symptoms.\par
\noindent\textit{Codes:} \texttt{chronic}.\quad \textit{Bearing:} (a) \textbf{N}, (b) \textbf{N}.\par
\noindent\textit{Rule trace:} (a) N: support {[}--{]}, challenge {[}--{]}; (b) N: motivation {[}--{]}, signal {[}--{]}, cost {[}--{]}; accommodated by Cor.~3.4 of EXI-10-a: {[}chronic{]}.\par}
\medskip
\needspace{6\baselineskip}
\noindent\textbf{Vod03}\enspace \citet{Vodanovich2003}.\par
{\small\noindent\textit{Design:} review.\quad \textit{Sample:} none (review).\quad \textit{Measure:} review of boredom scales.\par
\noindent\textit{Finding.} Reviews the psychometric properties of trait and state boredom measures and recommends further validation work.\par
\noindent\textit{Codes:} \texttt{meas}.\quad \textit{Bearing:} (a) \textbf{N}, (b) \textbf{N}.\par
\noindent\textit{Rule trace:} (a) N: support {[}--{]}, challenge {[}--{]}; (b) N: motivation {[}--{]}, signal {[}--{]}, cost {[}--{]}.\par}
\medskip
\needspace{6\baselineskip}
\noindent\textbf{Dan05}\enspace \citet{DanckertAllman2005}.\par
{\small\noindent\textit{Design:} correlational.\quad \textit{Sample:} healthy adults high or low in boredom proneness.\quad \textit{Measure:} time-estimation and attention tasks; boredom proneness.\par
\noindent\textit{Finding.} High and low boredom-prone participants did not differ in the temporal allocation of attention, but high boredom-prone participants made larger errors in estimating temporal intervals.\par
\noindent\textit{Codes:} \texttt{distinct}.\quad \textit{Bearing:} (a) \textbf{N}, (b) \textbf{S}.\par
\noindent\textit{Rule trace:} (a) N: support {[}--{]}, challenge {[}--{]}; (b) S: motivation {[}--{]}, signal {[}distinct{]}, cost {[}--{]}.\par}
\medskip
\needspace{6\baselineskip}
\noindent\textbf{Eas07}\enspace \citet{Eastwood2007}.\par
{\small\noindent\textit{Design:} correlational.\quad \textit{Sample:} undergraduates.\quad \textit{Measure:} BPS; alexithymia scale.\par
\noindent\textit{Finding.} Boredom proneness was associated with difficulty in identifying and describing one's feelings.\par
\noindent\textit{Codes:} \texttt{person}.\quad \textit{Bearing:} (a) \textbf{N}, (b) \textbf{N}.\par
\noindent\textit{Rule trace:} (a) N: support {[}--{]}, challenge {[}--{]}; (b) N: motivation {[}--{]}, signal {[}--{]}, cost {[}--{]}; accommodated by Cor.~3.4 of EXI-10-a: {[}person{]}.\par}
\medskip
\needspace{6\baselineskip}
\noindent\textbf{Eng08}\enspace \citet{EngeserRheinberg2008}.\par
{\small\noindent\textit{Design:} correlational.\quad \textit{Sample:} students and adults in learning and leisure settings.\quad \textit{Measure:} flow and challenge--skill ratings.\par
\noindent\textit{Finding.} Flow was related to the balance of challenge and skill, with the relation moderated by individual and situational factors.\par
\noindent\textit{Codes:} \texttt{fit-}.\quad \textit{Bearing:} (a) \textbf{S}, (b) \textbf{N}.\par
\noindent\textit{Rule trace:} (a) S: support {[}fit-{]}, challenge {[}--{]}; (b) N: motivation {[}--{]}, signal {[}--{]}, cost {[}--{]}.\par}
\medskip
\needspace{6\baselineskip}
\noindent\textbf{Fah09}\enspace \citet{Fahlman2009}.\par
{\small\noindent\textit{Design:} correlational.\quad \textit{Sample:} undergraduates (cross-sectional, longitudinal and experimental studies).\quad \textit{Measure:} BPS; meaning-in-life measures.\par
\noindent\textit{Finding.} Lack of meaning in life predicted boredom, including over time and, in an experimental study, when meaning was manipulated.\par
\noindent\textit{Codes:} \texttt{meaning+}.\quad \textit{Bearing:} (a) \textbf{S}, (b) \textbf{N}.\par
\noindent\textit{Rule trace:} (a) S: support {[}meaning+{]}, challenge {[}--{]}; (b) N: motivation {[}--{]}, signal {[}--{]}, cost {[}--{]}.\par}
\medskip
\needspace{6\baselineskip}
\noindent\textbf{Pek10}\enspace \citet{Pekrun2010}.\par
{\small\noindent\textit{Design:} correlational.\quad \textit{Sample:} university students.\quad \textit{Measure:} achievement-emotion questionnaire scales.\par
\noindent\textit{Finding.} Boredom in achievement settings was related to appraisals of control and of the value of the activity, and predicted lower achievement.\par
\noindent\textit{Codes:} \texttt{meaning+ perf-}.\quad \textit{Bearing:} (a) \textbf{S}, (b) \textbf{N}.\par
\noindent\textit{Rule trace:} (a) S: support {[}meaning+{]}, challenge {[}--{]}; (b) N: motivation {[}--{]}, signal {[}--{]}, cost {[}--{]}; accommodated by Cor.~3.4 of EXI-10-a: {[}perf-{]}.\par}
\medskip
\needspace{6\baselineskip}
\noindent\textbf{Gol11}\enspace \citet{Goldberg2011}.\par
{\small\noindent\textit{Design:} psychometric.\quad \textit{Sample:} undergraduates.\quad \textit{Measure:} boredom, apathy, anhedonia and depression scales.\par
\noindent\textit{Finding.} Boredom was related to, but separable from, apathy, anhedonia and depression.\par
\noindent\textit{Codes:} \texttt{distinct}.\quad \textit{Bearing:} (a) \textbf{N}, (b) \textbf{S}.\par
\noindent\textit{Rule trace:} (a) N: support {[}--{]}, challenge {[}--{]}; (b) S: motivation {[}--{]}, signal {[}distinct{]}, cost {[}--{]}.\par}
\medskip
\needspace{6\baselineskip}
\noindent\textbf{Kel11}\enspace \citet{Keller2011}.\par
{\small\noindent\textit{Design:} experiment.\quad \textit{Sample:} laboratory sample.\quad \textit{Measure:} flow ratings; heart-rate variability; salivary cortisol.\par
\noindent\textit{Finding.} A task whose demands were matched to the participant's skill, compared with too easy and too hard versions, produced flow and a distinct physiological profile.\par
\noindent\textit{Codes:} \texttt{fit-}.\quad \textit{Bearing:} (a) \textbf{S}, (b) \textbf{N}.\par
\noindent\textit{Rule trace:} (a) S: support {[}fit-{]}, challenge {[}--{]}; (b) N: motivation {[}--{]}, signal {[}--{]}, cost {[}--{]}.\par}
\medskip
\needspace{6\baselineskip}
\noindent\textbf{Til11}\enspace \citet{vanTilburgIgou2011}.\par
{\small\noindent\textit{Design:} experiment.\quad \textit{Sample:} student and adult samples.\quad \textit{Measure:} boredom inductions; in-group evaluation.\par
\noindent\textit{Finding.} Induced boredom increased the favoring of the in-group, interpreted as a way of re-establishing meaning.\par
\noindent\textit{Codes:} \texttt{meaningreg}.\quad \textit{Bearing:} (a) \textbf{N}, (b) \textbf{S}.\par
\noindent\textit{Rule trace:} (a) N: support {[}--{]}, challenge {[}--{]}; (b) S: motivation {[}meaningreg{]}, signal {[}--{]}, cost {[}--{]}.\par}
\medskip
\needspace{6\baselineskip}
\noindent\textbf{Eas12}\enspace \citet{EastwoodEtAl2012}.\par
{\small\noindent\textit{Design:} theory.\quad \textit{Sample:} none (theory).\quad \textit{Measure:} --.\par
\noindent\textit{Finding.} Defines boredom as the aversive experience of wanting, but being unable, to engage in satisfying activity, and grounds it in failures of attention.\par
\noindent\textit{Codes:} \texttt{attn signal}.\quad \textit{Bearing:} (a) \textbf{S}, (b) \textbf{S}.\par
\noindent\textit{Rule trace:} (a) S: support {[}attn{]}, challenge {[}--{]}; (b) S: motivation {[}--{]}, signal {[}signal{]}, cost {[}--{]}.\par}
\medskip
\needspace{6\baselineskip}
\noindent\textbf{Mal12}\enspace \citet{Malkovsky2012}.\par
{\small\noindent\textit{Design:} experiment.\quad \textit{Sample:} undergraduates.\quad \textit{Measure:} BPS; sustained-attention task.\par
\noindent\textit{Finding.} Boredom proneness was associated with failures of sustained attention.\par
\noindent\textit{Codes:} \texttt{attn person}.\quad \textit{Bearing:} (a) \textbf{S}, (b) \textbf{N}.\par
\noindent\textit{Rule trace:} (a) S: support {[}attn{]}, challenge {[}--{]}; (b) N: motivation {[}--{]}, signal {[}--{]}, cost {[}--{]}; accommodated by Cor.~3.4 of EXI-10-a: {[}person{]}.\par}
\medskip
\needspace{6\baselineskip}
\noindent\textbf{Til12}\enspace \citet{vanTilburgIgou2012}.\par
{\small\noindent\textit{Design:} experiment.\quad \textit{Sample:} student and adult samples.\quad \textit{Measure:} recalled and induced boredom; appraisal ratings.\par
\noindent\textit{Finding.} Boredom was characterized, relative to other negative emotions, by perceived lack of challenge and lack of meaning, and by a wish to do something else.\par
\noindent\textit{Codes:} \texttt{under+ meaning+ signal switch}.\quad \textit{Bearing:} (a) \textbf{S}, (b) \textbf{S}.\par
\noindent\textit{Rule trace:} (a) S: support {[}meaning+, under+{]}, challenge {[}--{]}; (b) S: motivation {[}switch{]}, signal {[}signal{]}, cost {[}--{]}.\par}
\medskip
\needspace{6\baselineskip}
\noindent\textbf{Ben13}\enspace \citet{BenchLench2013}.\par
{\small\noindent\textit{Design:} theory.\quad \textit{Sample:} none (theory and review).\quad \textit{Measure:} --.\par
\noindent\textit{Finding.} Argues that boredom motivates the pursuit of new goals when the current situation has ceased to be rewarding.\par
\noindent\textit{Codes:} \texttt{switch}.\quad \textit{Bearing:} (a) \textbf{N}, (b) \textbf{S}.\par
\noindent\textit{Rule trace:} (a) N: support {[}--{]}, challenge {[}--{]}; (b) S: motivation {[}switch{]}, signal {[}--{]}, cost {[}--{]}.\par}
\medskip
\needspace{6\baselineskip}
\noindent\textbf{Fah13}\enspace \citet{Fahlman2013}.\par
{\small\noindent\textit{Design:} psychometric.\quad \textit{Sample:} undergraduates and community adults.\quad \textit{Measure:} Multidimensional State Boredom Scale (MSBS).\par
\noindent\textit{Finding.} Develops a state boredom scale with factors for disengagement, high and low arousal, inattention and time perception.\par
\noindent\textit{Codes:} \texttt{meas}.\quad \textit{Bearing:} (a) \textbf{N}, (b) \textbf{N}.\par
\noindent\textit{Rule trace:} (a) N: support {[}--{]}, challenge {[}--{]}; (b) N: motivation {[}--{]}, signal {[}--{]}, cost {[}--{]}.\par}
\medskip
\needspace{6\baselineskip}
\noindent\textbf{Til13}\enspace \citet{vanTilburgIgouSedikides2013}.\par
{\small\noindent\textit{Design:} experiment.\quad \textit{Sample:} student and adult samples.\quad \textit{Measure:} boredom inductions; nostalgia; meaning.\par
\noindent\textit{Finding.} Induced boredom increased nostalgia, which in turn increased perceived meaning.\par
\noindent\textit{Codes:} \texttt{meaningreg}.\quad \textit{Bearing:} (a) \textbf{N}, (b) \textbf{S}.\par
\noindent\textit{Rule trace:} (a) N: support {[}--{]}, challenge {[}--{]}; (b) S: motivation {[}meaningreg{]}, signal {[}--{]}, cost {[}--{]}.\par}
\medskip
\needspace{6\baselineskip}
\noindent\textbf{Elp14}\enspace \citet{Elpidorou2014}.\par
{\small\noindent\textit{Design:} theory.\quad \textit{Sample:} none (theory).\quad \textit{Measure:} --.\par
\noindent\textit{Finding.} Argues that boredom is a regulatory state that informs the agent that her situation does not fit her goals and motivates a change.\par
\noindent\textit{Codes:} \texttt{signal switch}.\quad \textit{Bearing:} (a) \textbf{N}, (b) \textbf{S}.\par
\noindent\textit{Rule trace:} (a) N: support {[}--{]}, challenge {[}--{]}; (b) S: motivation {[}switch{]}, signal {[}signal{]}, cost {[}--{]}.\par}
\medskip
\needspace{6\baselineskip}
\noindent\textbf{Ger14}\enspace \citet{Gerritsen2014}.\par
{\small\noindent\textit{Design:} correlational.\quad \textit{Sample:} undergraduates.\quad \textit{Measure:} boredom proneness; inattention, hyperactivity, impulsivity and executive-function measures.\par
\noindent\textit{Finding.} Inattention, hyperactivity and executive dysfunction predicted boredom proneness, largely through the variance they share.\par
\noindent\textit{Codes:} \texttt{person}.\quad \textit{Bearing:} (a) \textbf{N}, (b) \textbf{N}.\par
\noindent\textit{Rule trace:} (a) N: support {[}--{]}, challenge {[}--{]}; (b) N: motivation {[}--{]}, signal {[}--{]}, cost {[}--{]}; accommodated by Cor.~3.4 of EXI-10-a: {[}person{]}.\par}
\medskip
\needspace{6\baselineskip}
\noindent\textbf{Goe14}\enspace \citet{Goetz2014}.\par
{\small\noindent\textit{Design:} experience sampling.\quad \textit{Sample:} high-school and university students.\quad \textit{Measure:} experience sampling of boredom, valence and arousal.\par
\noindent\textit{Finding.} Distinguishes several types of boredom differing in valence and arousal, among them a searching type in which the subject looks for alternatives.\par
\noindent\textit{Codes:} \texttt{meas seek}.\quad \textit{Bearing:} (a) \textbf{N}, (b) \textbf{S}.\par
\noindent\textit{Rule trace:} (a) N: support {[}--{]}, challenge {[}--{]}; (b) S: motivation {[}seek{]}, signal {[}--{]}, cost {[}--{]}.\par}
\medskip
\needspace{6\baselineskip}
\noindent\textbf{Man14}\enspace \citet{MannCadman2014}.\par
{\small\noindent\textit{Design:} experiment.\quad \textit{Sample:} adult participants.\quad \textit{Measure:} boredom induction; divergent-thinking task.\par
\noindent\textit{Finding.} Participants who first performed a boring task produced more ideas on a subsequent divergent-thinking task than controls.\par
\noindent\textit{Codes:} \texttt{creat}.\quad \textit{Bearing:} (a) \textbf{N}, (b) \textbf{S}.\par
\noindent\textit{Rule trace:} (a) N: support {[}--{]}, challenge {[}--{]}; (b) S: motivation {[}creat{]}, signal {[}--{]}, cost {[}--{]}.\par}
\medskip
\needspace{6\baselineskip}
\noindent\textbf{Mer14}\enspace \citet{MercerLynn2014}.\par
{\small\noindent\textit{Design:} correlational.\quad \textit{Sample:} undergraduates.\quad \textit{Measure:} trait and state boredom; situational ratings.\par
\noindent\textit{Finding.} Both trait boredom proneness and features of the situation contributed to state boredom.\par
\noindent\textit{Codes:} \texttt{relational person}.\quad \textit{Bearing:} (a) \textbf{S}, (b) \textbf{N}.\par
\noindent\textit{Rule trace:} (a) S: support {[}relational{]}, challenge {[}--{]}; (b) N: motivation {[}--{]}, signal {[}--{]}, cost {[}--{]}; accommodated by Cor.~3.4 of EXI-10-a: {[}person{]}.\par}
\medskip
\needspace{6\baselineskip}
\noindent\textbf{Mrf14}\enspace \citet{MerrifieldDanckert2014}.\par
{\small\noindent\textit{Design:} physiological / imaging.\quad \textit{Sample:} undergraduates.\quad \textit{Measure:} physiological measures during mood inductions.\par
\noindent\textit{Finding.} Induced boredom produced a physiological profile that differed from that of induced sadness.\par
\noindent\textit{Codes:} \texttt{distinct}.\quad \textit{Bearing:} (a) \textbf{N}, (b) \textbf{S}.\par
\noindent\textit{Rule trace:} (a) N: support {[}--{]}, challenge {[}--{]}; (b) S: motivation {[}--{]}, signal {[}distinct{]}, cost {[}--{]}.\par}
\medskip
\needspace{6\baselineskip}
\noindent\textbf{Wil14}\enspace \citet{WilsonEtAl2014}.\par
{\small\noindent\textit{Design:} experiment.\quad \textit{Sample:} undergraduates and community adults.\quad \textit{Measure:} enjoyment ratings; self-administered electric shocks.\par
\noindent\textit{Finding.} Participants generally did not enjoy spending several minutes alone with their thoughts, and many chose to give themselves mild electric shocks rather than sit without stimulation.\par
\noindent\textit{Codes:} \texttt{seek harm}.\quad \textit{Bearing:} (a) \textbf{N}, (b) \textbf{S/C}.\par
\noindent\textit{Rule trace:} (a) N: support {[}--{]}, challenge {[}--{]}; (b) S/C: motivation {[}seek{]}, signal {[}--{]}, cost {[}harm{]}.\par}
\medskip
\needspace{6\baselineskip}
\noindent\textbf{Hav15}\enspace \citet{Havermans2015}.\par
{\small\noindent\textit{Design:} experiment.\quad \textit{Sample:} laboratory sample.\quad \textit{Measure:} boredom induction; food intake; self-administered shocks.\par
\noindent\textit{Finding.} Under a boredom induction, participants ate more and administered more shocks to themselves than in a comparison condition.\par
\noindent\textit{Codes:} \texttt{seek harm}.\quad \textit{Bearing:} (a) \textbf{N}, (b) \textbf{S/C}.\par
\noindent\textit{Rule trace:} (a) N: support {[}--{]}, challenge {[}--{]}; (b) S/C: motivation {[}seek{]}, signal {[}--{]}, cost {[}harm{]}.\par}
\medskip
\needspace{6\baselineskip}
\noindent\textbf{Ned16}\enspace \citet{Nederkoorn2016}.\par
{\small\noindent\textit{Design:} experiment.\quad \textit{Sample:} laboratory sample.\quad \textit{Measure:} boring versus comparison film; self-administered shocks.\par
\noindent\textit{Finding.} A boring film led to more self-administered painful stimulation than a comparison film.\par
\noindent\textit{Codes:} \texttt{seek harm}.\quad \textit{Bearing:} (a) \textbf{N}, (b) \textbf{S/C}.\par
\noindent\textit{Rule trace:} (a) N: support {[}--{]}, challenge {[}--{]}; (b) S/C: motivation {[}seek{]}, signal {[}--{]}, cost {[}harm{]}.\par}
\medskip
\needspace{6\baselineskip}
\noindent\textbf{Str16}\enspace \citet{Struk2016}.\par
{\small\noindent\textit{Design:} correlational.\quad \textit{Sample:} undergraduates.\quad \textit{Measure:} BPS; regulatory-mode scales.\par
\noindent\textit{Finding.} Boredom proneness was associated with a regulatory profile of low locomotion and high assessment.\par
\noindent\textit{Codes:} \texttt{person}.\quad \textit{Bearing:} (a) \textbf{N}, (b) \textbf{N}.\par
\noindent\textit{Rule trace:} (a) N: support {[}--{]}, challenge {[}--{]}; (b) N: motivation {[}--{]}, signal {[}--{]}, cost {[}--{]}; accommodated by Cor.~3.4 of EXI-10-a: {[}person{]}.\par}
\medskip
\needspace{6\baselineskip}
\noindent\textbf{Til16}\enspace \citet{vanTilburgIgou2016}.\par
{\small\noindent\textit{Design:} experiment.\quad \textit{Sample:} student and adult samples.\quad \textit{Measure:} boredom inductions; political orientation.\par
\noindent\textit{Finding.} Boredom was associated with more extreme political orientations, interpreted as a way of re-establishing meaning.\par
\noindent\textit{Codes:} \texttt{meaningreg}.\quad \textit{Bearing:} (a) \textbf{N}, (b) \textbf{S}.\par
\noindent\textit{Rule trace:} (a) N: support {[}--{]}, challenge {[}--{]}; (b) S: motivation {[}meaningreg{]}, signal {[}--{]}, cost {[}--{]}.\par}
\medskip
\needspace{6\baselineskip}
\noindent\textbf{Tze16}\enspace \citet{Tze2016}.\par
{\small\noindent\textit{Design:} meta-analysis.\quad \textit{Sample:} school and university students (pooled).\quad \textit{Measure:} boredom and academic outcomes.\par
\noindent\textit{Finding.} Across studies, boredom was negatively related to academic outcomes.\par
\noindent\textit{Codes:} \texttt{perf-}.\quad \textit{Bearing:} (a) \textbf{N}, (b) \textbf{N}.\par
\noindent\textit{Rule trace:} (a) N: support {[}--{]}, challenge {[}--{]}; (b) N: motivation {[}--{]}, signal {[}--{]}, cost {[}--{]}; accommodated by Cor.~3.4 of EXI-10-a: {[}perf-{]}.\par}
\medskip
\needspace{6\baselineskip}
\noindent\textbf{Vod16}\enspace \citet{VodanovichWatt2016}.\par
{\small\noindent\textit{Design:} review.\quad \textit{Sample:} none (review).\quad \textit{Measure:} review of self-report boredom measures.\par
\noindent\textit{Finding.} Updates the review of trait and state boredom scales, including newer instruments.\par
\noindent\textit{Codes:} \texttt{meas}.\quad \textit{Bearing:} (a) \textbf{N}, (b) \textbf{N}.\par
\noindent\textit{Rule trace:} (a) N: support {[}--{]}, challenge {[}--{]}; (b) N: motivation {[}--{]}, signal {[}--{]}, cost {[}--{]}.\par}
\medskip
\needspace{6\baselineskip}
\noindent\textbf{Chi17}\enspace \citet{Chin2017}.\par
{\small\noindent\textit{Design:} experience sampling.\quad \textit{Sample:} adults (smartphone sampling).\quad \textit{Measure:} momentary boredom reports.\par
\noindent\textit{Finding.} Boredom was common in everyday life and varied with activity and social context.\par
\noindent\textit{Codes:} \texttt{relational}.\quad \textit{Bearing:} (a) \textbf{S}, (b) \textbf{N}.\par
\noindent\textit{Rule trace:} (a) S: support {[}relational{]}, challenge {[}--{]}; (b) N: motivation {[}--{]}, signal {[}--{]}, cost {[}--{]}.\par}
\medskip
\needspace{6\baselineskip}
\noindent\textbf{Isa17}\enspace \citet{Isacescu2017}.\par
{\small\noindent\textit{Design:} correlational.\quad \textit{Sample:} undergraduates.\quad \textit{Measure:} boredom proneness; self-control and affect measures.\par
\noindent\textit{Finding.} Self-control and affective variables predicted boredom proneness.\par
\noindent\textit{Codes:} \texttt{person}.\quad \textit{Bearing:} (a) \textbf{N}, (b) \textbf{N}.\par
\noindent\textit{Rule trace:} (a) N: support {[}--{]}, challenge {[}--{]}; (b) N: motivation {[}--{]}, signal {[}--{]}, cost {[}--{]}; accommodated by Cor.~3.4 of EXI-10-a: {[}person{]}.\par}
\medskip
\needspace{6\baselineskip}
\noindent\textbf{Str17}\enspace \citet{Struk2017}.\par
{\small\noindent\textit{Design:} psychometric.\quad \textit{Sample:} undergraduates and online adults.\quad \textit{Measure:} Short Boredom Proneness Scale (SBPS).\par
\noindent\textit{Finding.} Develops an eight-item unidimensional short form of the boredom proneness construct with good psychometric properties.\par
\noindent\textit{Codes:} \texttt{meas}.\quad \textit{Bearing:} (a) \textbf{N}, (b) \textbf{N}.\par
\noindent\textit{Rule trace:} (a) N: support {[}--{]}, challenge {[}--{]}; (b) N: motivation {[}--{]}, signal {[}--{]}, cost {[}--{]}.\par}
\medskip
\needspace{6\baselineskip}
\noindent\textbf{Til17}\enspace \citet{vanTilburgIgou2017}.\par
{\small\noindent\textit{Design:} experiment.\quad \textit{Sample:} student and adult samples.\quad \textit{Measure:} recalled emotions; appraisal and motivation ratings.\par
\noindent\textit{Finding.} Boredom differed from other negative emotions in its profile, notably in perceived meaninglessness and in the wish to pursue a different goal.\par
\noindent\textit{Codes:} \texttt{meaning+ signal switch}.\quad \textit{Bearing:} (a) \textbf{S}, (b) \textbf{S}.\par
\noindent\textit{Rule trace:} (a) S: support {[}meaning+{]}, challenge {[}--{]}; (b) S: motivation {[}switch{]}, signal {[}signal{]}, cost {[}--{]}.\par}
\medskip
\needspace{6\baselineskip}
\noindent\textbf{DnH18}\enspace \citet{Danckert2018restless}.\par
{\small\noindent\textit{Design:} correlational.\quad \textit{Sample:} undergraduates and online adults.\quad \textit{Measure:} state boredom; arousal and restlessness reports.\par
\noindent\textit{Finding.} Boredom was associated with low arousal combined with restlessness.\par
\noindent\textit{Codes:} \texttt{distinct}.\quad \textit{Bearing:} (a) \textbf{N}, (b) \textbf{S}.\par
\noindent\textit{Rule trace:} (a) N: support {[}--{]}, challenge {[}--{]}; (b) S: motivation {[}--{]}, signal {[}distinct{]}, cost {[}--{]}.\par}
\medskip
\needspace{6\baselineskip}
\noindent\textbf{DnM18}\enspace \citet{DanckertMerrifield2018}.\par
{\small\noindent\textit{Design:} physiological / imaging.\quad \textit{Sample:} adult participants.\quad \textit{Measure:} fMRI during boredom induction.\par
\noindent\textit{Finding.} Boredom induction was associated with activity in the default mode network.\par
\noindent\textit{Codes:} \texttt{neuro}.\quad \textit{Bearing:} (a) \textbf{N}, (b) \textbf{N}.\par
\noindent\textit{Rule trace:} (a) N: support {[}--{]}, challenge {[}--{]}; (b) N: motivation {[}--{]}, signal {[}--{]}, cost {[}--{]}.\par}
\medskip
\needspace{6\baselineskip}
\noindent\textbf{Dnk18}\enspace \citet{DanckertMugon2018}.\par
{\small\noindent\textit{Design:} theory.\quad \textit{Sample:} none (theory).\quad \textit{Measure:} --.\par
\noindent\textit{Finding.} Argues that boredom signals that one's current engagement is unsatisfactory and prompts the exploration of alternatives.\par
\noindent\textit{Codes:} \texttt{signal seek}.\quad \textit{Bearing:} (a) \textbf{N}, (b) \textbf{S}.\par
\noindent\textit{Rule trace:} (a) N: support {[}--{]}, challenge {[}--{]}; (b) S: motivation {[}seek{]}, signal {[}signal{]}, cost {[}--{]}.\par}
\medskip
\needspace{6\baselineskip}
\noindent\textbf{Elp18}\enspace \citet{Elpidorou2018}.\par
{\small\noindent\textit{Design:} theory.\quad \textit{Sample:} none (theory).\quad \textit{Measure:} --.\par
\noindent\textit{Finding.} Argues that boredom is valuable as a signal of a mismatch between the agent's situation and her desires, and as a motivation to change it.\par
\noindent\textit{Codes:} \texttt{signal switch}.\quad \textit{Bearing:} (a) \textbf{N}, (b) \textbf{S}.\par
\noindent\textit{Rule trace:} (a) N: support {[}--{]}, challenge {[}--{]}; (b) S: motivation {[}switch{]}, signal {[}signal{]}, cost {[}--{]}.\par}
\medskip
\needspace{6\baselineskip}
\noindent\textbf{Hun18}\enspace \citet{HunterEastwood2018}.\par
{\small\noindent\textit{Design:} experiment.\quad \textit{Sample:} undergraduates.\quad \textit{Measure:} trait and state boredom; sustained-attention task.\par
\noindent\textit{Finding.} Examines the relations of trait and state boredom to sustained attention; attentional performance and boredom were related.\par
\noindent\textit{Codes:} \texttt{attn}.\quad \textit{Bearing:} (a) \textbf{S}, (b) \textbf{N}.\par
\noindent\textit{Rule trace:} (a) S: support {[}attn{]}, challenge {[}--{]}; (b) N: motivation {[}--{]}, signal {[}--{]}, cost {[}--{]}.\par}
\medskip
\needspace{6\baselineskip}
\noindent\textbf{Mug18}\enspace \citet{Mugon2018}.\par
{\small\noindent\textit{Design:} correlational.\quad \textit{Sample:} undergraduates.\quad \textit{Measure:} BPS; regulatory modes.\par
\noindent\textit{Finding.} Boredom proneness was associated with difficulty in initiating action (low locomotion).\par
\noindent\textit{Codes:} \texttt{person}.\quad \textit{Bearing:} (a) \textbf{N}, (b) \textbf{N}.\par
\noindent\textit{Rule trace:} (a) N: support {[}--{]}, challenge {[}--{]}; (b) N: motivation {[}--{]}, signal {[}--{]}, cost {[}--{]}; accommodated by Cor.~3.4 of EXI-10-a: {[}person{]}.\par}
\medskip
\needspace{6\baselineskip}
\noindent\textbf{Raf18}\enspace \citet{Raffaelli2018}.\par
{\small\noindent\textit{Design:} review.\quad \textit{Sample:} none (review).\quad \textit{Measure:} --.\par
\noindent\textit{Finding.} Reviews what is and is not known about boredom across definitional, measurement and neuroscientific work.\par
\noindent\textit{Codes:} \texttt{meas}.\quad \textit{Bearing:} (a) \textbf{N}, (b) \textbf{N}.\par
\noindent\textit{Rule trace:} (a) N: support {[}--{]}, challenge {[}--{]}; (b) N: motivation {[}--{]}, signal {[}--{]}, cost {[}--{]}.\par}
\medskip
\needspace{6\baselineskip}
\noindent\textbf{Wes18}\enspace \citet{WestgateWilson2018}.\par
{\small\noindent\textit{Design:} theory.\quad \textit{Sample:} none (theory).\quad \textit{Measure:} --.\par
\noindent\textit{Finding.} The MAC model: boredom arises from a mismatch between cognitive demands and available resources (too low or too high) and from a lack of meaning, both signaling attentional or value deficits.\par
\noindent\textit{Codes:} \texttt{under+ over+ meaning+ attn signal}.\quad \textit{Bearing:} (a) \textbf{S}, (b) \textbf{S}.\par
\noindent\textit{Rule trace:} (a) S: support {[}attn, meaning+, over+, under+{]}, challenge {[}--{]}; (b) S: motivation {[}--{]}, signal {[}signal{]}, cost {[}--{]}.\par}
\medskip
\needspace{6\baselineskip}
\noindent\textbf{Ben19}\enspace \citet{BenchLench2019}.\par
{\small\noindent\textit{Design:} experiment.\quad \textit{Sample:} undergraduates.\quad \textit{Measure:} boredom induction; choice of novel stimuli.\par
\noindent\textit{Finding.} Induced boredom increased the choice of novel experiences, even unpleasant ones.\par
\noindent\textit{Codes:} \texttt{seek}.\quad \textit{Bearing:} (a) \textbf{N}, (b) \textbf{S}.\par
\noindent\textit{Rule trace:} (a) N: support {[}--{]}, challenge {[}--{]}; (b) S: motivation {[}seek{]}, signal {[}--{]}, cost {[}--{]}.\par}
\medskip
\needspace{6\baselineskip}
\noindent\textbf{DnE20}\enspace \citet{DanckertEastwood2020}.\par
{\small\noindent\textit{Design:} review.\quad \textit{Sample:} none (review).\quad \textit{Measure:} --.\par
\noindent\textit{Finding.} Synthesizes research on boredom as a call to action: a signal that one is not engaged that prompts the search for engagement.\par
\noindent\textit{Codes:} \texttt{attn signal seek}.\quad \textit{Bearing:} (a) \textbf{S}, (b) \textbf{S}.\par
\noindent\textit{Rule trace:} (a) S: support {[}attn{]}, challenge {[}--{]}; (b) S: motivation {[}seek{]}, signal {[}signal{]}, cost {[}--{]}.\par}
\medskip
\needspace{6\baselineskip}
\noindent\textbf{Wol20}\enspace \citet{Wolff2020}.\par
{\small\noindent\textit{Design:} correlational.\quad \textit{Sample:} adults (online).\quad \textit{Measure:} boredom proneness; self-control; adherence to distancing rules.\par
\noindent\textit{Finding.} Higher boredom proneness and lower self-control were associated with weaker adherence to social-distancing guidelines.\par
\noindent\textit{Codes:} \texttt{chronic}.\quad \textit{Bearing:} (a) \textbf{N}, (b) \textbf{N}.\par
\noindent\textit{Rule trace:} (a) N: support {[}--{]}, challenge {[}--{]}; (b) N: motivation {[}--{]}, signal {[}--{]}, cost {[}--{]}; accommodated by Cor.~3.4 of EXI-10-a: {[}chronic{]}.\par}
\medskip
\needspace{6\baselineskip}
\noindent\textbf{Boy21}\enspace \citet{Boylan2021}.\par
{\small\noindent\textit{Design:} correlational.\quad \textit{Sample:} adults (online).\quad \textit{Measure:} boredom proneness; desire to act; rule-breaking.\par
\noindent\textit{Finding.} Trait boredom proneness was associated with a desire to act and with self-reported breaking of pandemic rules.\par
\noindent\textit{Codes:} \texttt{chronic seek}.\quad \textit{Bearing:} (a) \textbf{N}, (b) \textbf{S}.\par
\noindent\textit{Rule trace:} (a) N: support {[}--{]}, challenge {[}--{]}; (b) S: motivation {[}seek{]}, signal {[}--{]}, cost {[}--{]}; accommodated by Cor.~3.4 of EXI-10-a: {[}chronic{]}.\par}
\medskip
\needspace{6\baselineskip}
\noindent\textbf{Tam21}\enspace \citet{Tam2021}.\par
{\small\noindent\textit{Design:} theory.\quad \textit{Sample:} none (theory).\quad \textit{Measure:} --.\par
\noindent\textit{Finding.} The Boredom Feedback Model: boredom is feedback that attention is not adequately engaged, prompting attempts to re-engage.\par
\noindent\textit{Codes:} \texttt{attn signal seek}.\quad \textit{Bearing:} (a) \textbf{S}, (b) \textbf{S}.\par
\noindent\textit{Rule trace:} (a) S: support {[}attn{]}, challenge {[}--{]}; (b) S: motivation {[}seek{]}, signal {[}signal{]}, cost {[}--{]}.\par}
\medskip
